\documentclass{article}
\usepackage{amsmath, amssymb, amsthm}
\usepackage{hyperref}
\usepackage{geometry}
\geometry{margin=1in}

\title{Erd\H{o}s Problem \#683}
\date{Last updated: 2025-09-04}
\author{Source: erdosproblems.com}

\begin{document}
\maketitle

\noindent\textbf{Status:} OPEN \quad \textbf{No prize}

\noindent\textbf{Tags:} number theory, primes, binomial coefficients

\medskip

\noindent\textbf{Problem Statement:}

Is it true that for every $1\leq k\leq n$ the largest prime divisor of $\binom{n}{k}$, say $P(\binom{n}{k})$, satisfies\[P\left(\binom{n}{k}\right)\geq \min(n-k+1, k^{1+c})\]for some constant $c&#62;0$?




A theorem of Sylvester and Schur (see \cite{Er34}) states that $P(\binom{n}{k})&gt;k$ if $k\leq n/2$. Erd\H{o}s \cite{Er55d} proved that there exists some $c&gt;0$ such that, whenever $k\leq n/2$,\[P\left(\binom{n}{k}\right)\gg k\log k.\]Erd\H{o}s \cite{Er79d} writes it 'seems certain' that this holds for every $c&gt;0$, with only a finite number of exceptions (depending on $c$). Standard heuristics on prime gaps suggest that the largest prime divisor of $\binom{n}{k}$ is, for $k\leq n/2$, in fact\[&gt;e^{c\sqrt{k}}\]for some constant $c&gt;0$. This is essentially equivalent to [961] .




References

[Er34] Erd\H{o}s, Paul, A Theorem of Sylvester and Schur . J. London Math. Soc. (1934), 282--288.

[Er55d] Erd\H{o}s, P., On consecutive integers . Nieuw Arch. Wisk. (3) (1955), 124--128.

[Er79d] Erd\H{o}s, P., Some unconventional problems in number theory . Acta Math. Acad. Sci. Hungar. (1979), 71-80.

\noindent\textbf{Related OEIS sequences:} A006530, A074399, A121359, possible


\bigskip
\noindent\small{Source: \url{https://www.erdosproblems.com/683}}
\end{document}
